\section{Three principles}\label{sec:three-principles}

After the symplectic primitive, scalar quotient, and residue
normalization are fixed, the local model is organized by three source
principles.  Each is physical in origin, but each is used here only
after translation into an algebraic hypothesis or target.  Together they
specify the hypotheses and algebraic targets for the open brane theory,
the closed Hamiltonian BF theory, the closed-open comparison, and the
cotangent coefficient module proved below.

\subsection{Dirac: a brane stack is a constrained probe}

Fix the formal holomorphic symplectic disk
\[
  \widehat{\C^2}_0=\operatorname{Spf}\C[[z_1,z_2]],
  \qquad \omega=dz_1\wedge dz_2.
\]
A stack of \(N\) branes supported on the topological line probes this
disk by replacing the normal coordinates by matrices
\[
  z_1\rightsquigarrow\phi_1,\qquad
  z_2\rightsquigarrow\phi_2,\qquad
  \phi_i\in\mathfrak{gl}_N .
\]
The pulled-back symplectic form is
\[
  \Omega_N=\operatorname{Tr}(d\phi_1\wedge d\phi_2),
\]
with primitive \(\operatorname{Tr}(\phi_1\,d\phi_2)\). Change of basis
in the probe stack is gauge. Its moment map is
\[
  \mu(\epsilon)=-\operatorname{Tr}(\epsilon[\phi_1,\phi_2]),
  \qquad \epsilon\in\mathfrak{gl}_N .
\]
Dirac reduction therefore imposes the first-class constraint
\[
  [\phi_1,\phi_2]=0.
\]
After the primitive \(\operatorname{Tr}(\phi_1\,d\phi_2)\) and moment
map convention have been fixed, the corresponding first-order
world-line action is
\[
  S_\partial
  =\int_\R\operatorname{Tr}
    \bigl(\phi_1\,d\phi_2+A[\phi_1,\phi_2]\bigr).
\]
In BV form the constraint is the Koszul differential
\[
  Q\psi=[\phi_1,\phi_2],\qquad Q\phi_1=Q\phi_2=0,
\]
where \(\psi=A^*\). Thus the open algebra is not postulated as a
trace algebra. It is the algebra of gauge-invariant functions on the
derived commuting-pair locus. Lemma~\ref{lem:dirac-probe-reduction}
and Proposition~\ref{prop:open-bv-truncation} give the formal proof.

\subsection{Witten convention: closed fields target central operations}

Witten's open string field theory supplies the cyclic open-field-theory
source convention, not a theorem imported into this local model.  The
closed-open expectation is that a closed insertion approaching a brane
does not merely evaluate to a boundary function; it should give an
operation on compatible boundary observables.  In a topological
direction the algebraic target for such an operation is centrality up to
derived Poisson homotopies.  The target tested in this paper is therefore
\[
  Z^{P_0}_{\mathrm{red,Ham}}(\mathrm{Obs}_{\mathrm{open}}),
\]
not \(\mathrm{Obs}_{\mathrm{open}}\) itself.  The centrality statement
proved below is the reduced CE/PV calculation in this formal Hamiltonian
sector; no global open string field theory or physical centrality theorem
is imported from Witten.

For the formal symplectic disk the infinitesimal closed symmetries are
Hamiltonian vector fields. Constants generate the zero vector field,
so the Hamiltonian Lie algebra is
\[
  \mathfrak h=\C[[z_1,z_2]]/\C\cdot1,
  \qquad
  \{f,g\}
  =\partial_{z_1}f\,\partial_{z_2}g
   -\partial_{z_2}f\,\partial_{z_1}g .
\]
In the local model the closed BV algebra of canonical transformations is
the shifted cotangent theory
\[
  \mathfrak g
  =\mathfrak h\ltimes\mathfrak h^\vee_{\mathrm{cont}}[1].
\]
The local theorem to be proved is the cochain identity
\[
  C^\bullet_{\mathrm{CE},\mathrm{cont}}(\mathfrak g)
  \cong
  \mathrm{PV}\bigl(\widehat{\mathbf S}_{\mathrm{Tate}}(\mathfrak h)\bigr)
  \simeq
  Z^{P_0}_{\mathrm{red,Ham}}
  \bigl(\widehat{\mathbf S}_{\mathrm{Tate}}(\mathfrak h)\bigr),
\]
with generator dictionary
\[
  c^I\mapsto\theta^I,\qquad u_I\mapsto O_I .
\]
The boundary Hamiltonian observable is the image of the shifted
cotangent coordinate \(u_I\). It is not the image of the Hamiltonian
ghost coordinate \(c^I\). This distinction is the type check behind
the factorization-current theorem.

The same type check blocks the unreduced lift.  A theorem at the
unreduced BV level would have to produce actual boundary representatives
for the shifted cotangent generators and homotopies making them central
against arbitrary open observables before primitive trace projection.
The reduced theorem below does not do this: it identifies the CE/PV
target and the reduced \(P_0\) current target, while the unreduced
centrality and QME compatibility remain separate obligations.

There is a sharper obstruction than a wrong label map. In the
polynomial unreduced BV/Koszul category the linear Hamiltonians act by
constant translations of \(\phi_2\) and \(\phi_1\), hence locally
nilpotently on every finite polynomial observable and on its
\(Q\)-cohomology. The principal-part cotangent module has the opposite
behavior:
\[
  z_1^r\cdot\rho_{a,b}
  =(-1)^r(b+1)\cdots(b+r)\rho_{a,b+r}\neq0,
  \qquad r\geq1,
\]
and similarly for \(z_2\). Thus a same-action unreduced lift cannot be
hidden inside polynomial one-\(\psi\) representatives. A corrected
boundary target must contain genuinely distributional principal-part
fields and must construct new centrality homotopies there.

The QME obstruction has an independent scalar component. With the
Capelli/Weyl convention used below,
\[
  \operatorname{Tr}(A\,XY)-\operatorname{Tr}(A\,YX)
  =\hbar N\,\operatorname{Tr}(A)
  \quad\bmod\ \mathcal W_N\widehat\mu(\lie{gl}_N).
\]
After inserting the central boundary ghost
\(\gamma_{\mathbf 1}\), \(|\gamma_{\mathbf 1}|=1\), the scalar contact
contributes
\[
  \operatorname{Ob}_{\mathrm{sc}}
  =
  \hbar N\,\omega(f,g)
  \int_\R a(t)b(t)\gamma_{\mathbf 1}(t)\,dt .
\]
Its associated graded class is \(\hbar N[\bar c]\). The reduced
Hamiltonian source has no constant Hamiltonian line on which the
primitive \(\eta(f)=-[f]_0\) could live. Known candidate mechanisms for
an ordinary \(\lie{gl}_N\) mixed brane-defect QME proof include a
closed-exchange term with associated graded \(-\hbar N[\bar c]\), a
supertrace or central-character normalization, or an unreduced scalar
primitive. The present manuscript records the class; it does not prove
that this list is exhaustive or that such a cancellation has been
constructed.

\subsection{Grothendieck: cotangent modes are residue currents}

The continuous dual of functions on a formal disk is not another copy
of polynomial functions. It is the local-cohomology module of residue
currents. For
\[
  R=\C[[z_1,z_2]],\qquad \mathfrak m=(z_1,z_2),
\]
the Matlis dualizing module is
\[
  H^2_{\mathfrak m}(R)\,dz_1dz_2
  =
  \bigoplus_{a,b\geq0}
  \C\,z_1^{-a-1}z_2^{-b-1}\,dz_1dz_2 .
\]
After removing the scalar Hamiltonian, the cotangent module is
\[
  \mathcal P
  =
  \operatorname{Ann}_{\mathrm{res}}(\C\cdot1)
  =
  \bigoplus_{a+b>0}\C\,\rho_{a,b},
  \qquad
  \rho_{a,b}=z_1^{-a-1}z_2^{-b-1}\,dz_1dz_2 .
\]
The residue pairing
\[
  \langle\rho_{a,b},z_1^m z_2^n\rangle
  =
  \operatorname{Res}_{0}
  z_1^{m-a-1}z_2^{n-b-1}\,dz_1dz_2
  =
  \delta_{a,m}\delta_{b,n}
\]
identifies \(\mathfrak h^\vee_{\mathrm{cont}}\) with \(\mathcal P\).
The coadjoint action is the lowering formula
\[
  z_1^p z_2^q\cdot\rho_{a,b}
  =
  -(pb-qa+p-q)\rho_{a-p+1,b-q+1}.
\]
This is the cotangent sector used by the closed theory. The
one-\(\psi\) open classes \(\Psi_{a,b}\) carry the same labels but a
different \(\mathfrak h\)-module structure. They form a PBW
special-fibre shadow; they do not realize the principal-part module.
The no-go theorem
\ref{thm:no-polynomial-realization-categorical} records the formal
obstruction.

\subsection{Deduction of the local package}

The three principles have disjoint jobs.

\begin{enumerate}
  \item Dirac determines the open BV complex and the stable trace
  observables.
  \item Witten's source convention motivates the closed-open target as a
  derived \(P_0\)-centre; the local CE/PV theorem supplies the proof in
  the reduced Hamiltonian sector.
  \item Grothendieck determines the continuous cotangent coefficients
  and separates principal parts from polynomial descendants.
\end{enumerate}

The rest of the paper proves the resulting package in independent
lanes: the Dirac probe theorem; the stable trace and one-\(\psi\)
calculation; the CE/PV derived-centre theorem; the Matlis
principal-part theorem; the brane-line factorization-current theorem;
the reduced principal-part defect-current theorem; the degree-zero
Moyal/Weyl quantum theorem; and the scalar \(U(1)\) anomaly theorem.
Open analytic obligations are named separately: the unreduced
cotangent factorization-centre morphism, the one-antifield quantum
lift, and the Costello brane-defect QME realization.
